\subsection{People Have a Right to Consume AND a Right to Produce}
\label{sec:right-consume}
The competing floor is a basic income, and the choice between them is usually argued on forecasts about how much work machines will take. It does not have to be, and the forecast is not the reason to prefer one.

A jobs guarantee is a public option. Medicare for all is one, and so is a public school: a thing the state provides directly, to anyone, at a standard it is answerable for, sitting in the same market as the private version and setting the floor that version has to beat. A basic income is not that. It is a transfer, and what it buys is whatever is for sale. Consuming public resources is not the same as being public — a charter school runs on public money and is not a public school — and a transfer spent at a private employer's store does not make the store public.

What employment carries besides money is the rest of it. A person in a guaranteed job is attached to other people by the work: colleagues, a schedule, a place, something produced that somebody needed. A transfer supplies the purchasing power and none of that, and the purchasing power was never the part in doubt.

A transfer also leaves a person outside the one body of law that reaches this. Anti-discrimination duties attach to employment. They do not attach to income. A guaranteed job puts somebody inside that architecture; a payment leaves them in a market, dealing with whoever will sell to them, on whatever terms. The architecture is contested and full of holes — it reaches employers above a size, and what it demands of a seller has always been narrower than what it demands of an employer — but a hole in a protection is something to litigate. A transfer does not itself provide the employment relationship and its associated protections. An income entitlement can still provide a claim to eligibility and payment. My argument for a guarantee is the additional right to produce, and the comparison of safeguards must state what each arrangement actually supplies.

Both proposals are also conditional, differently. A basic income is conditional on continued political willingness to fund it, which the aggregation machinery can withdraw in a budget; a jobs guarantee is conditional on a placement decision, which is a lever, and a dangerous one under a hostile administration, but a lever that leaves a decision-maker, a record, a counterparty and an appeal behind it. A basic income can be entrenched too, and a constitutional one would be harder to withdraw than a statutory guarantee. Compare the two designs at both points of withdrawal: who can end the funding, and who can deny one person access. A basic income goes wholesale, in a budget line, from everybody at once, and a placement is refused to a person, by somebody, on a date, which may expose an official, a date, and an appeal. The comparison has to show what the income arrangement exposes at the corresponding points.

A decision someone can contest gives dependence a point of resistance. In the comparison made here, the advantage belongs to the arrangement that leaves a usable record and someone who must answer for it. It is also, for the same reason, the more demanding proposal to build: everything in the comparison rests on the appeal being real, and the previous section says what makes it real.

That standard now has a court behind it, and on the half of the case that travels. When the Supreme Court kept a Federal Reserve governor in her post in 2026 against a removal for cause, what it held she was owed was not the post but the artifact: “some explanation of the evidence at issue, some avenue for a response, and a deadline by which a response would be due” \autocite{cook2026}. What does not travel is this: the Fed's protection rests on a pedigree no body chartered later can acquire. What a removal has to produce before it counts is not tied to that pedigree.

The argument concerns a right to produce, which a transfer does not supply and a guarantee does. A basic income is the right to consume, and the right to consume is increasingly a credit at somebody's store.

I put a jobs guarantee ahead of a basic income on those grounds. One remaining measure does work the floor cannot: broadened social insurance — unemployment benefits, healthcare, housing — reaches the people who cannot work or choose not to, whom a guarantee by construction does not reach at all.
